\documentclass[11pt]{article}
\usepackage{icmlko}
\begin{document}

\title{시계 둘이 거꾸로 간다 --- 어제 채택한 것을 오늘 다시 골랐다면 안 골랐을 것이다}
\author{월드모델 실험실}
\date{노트 500 · 2026년 8월 3일}
\maketitle

\begin{abstract}
노트 499 가 남긴 마지막 물음 --- \emph{안쪽 검증 시기를 옮기면 다른 표가
나오나.} 같은 후보 $21$ 짝을 안쪽 검증 $2022\!\sim\!2024$(안쪽 학습
$<\!2022$)로 다시 쟀다. \textbf{유보는 전혀 안 봤다.}

\textbf{두 시기의 순위 스피어만이 $\mathbf{-0.320}$ 이다.} 불안정한 게
아니라 \emph{거꾸로} 간다.

\begin{center}
\begin{tabular}{lrr}
\hline
(도메인, 축) & $2023\!\sim\!2025$ & $2022\!\sim\!2024$ \\
\hline
\textbf{애니 트렌드} & $\mathbf{+0.0042}$ ($1$ 등) & $\mathbf{-0.0064}$ ($21$ 등) \\
모바일 media\_push & $-0.0016$ & $+0.0035$ ($1$ 등) \\
모바일 태그 & $+0.0011$ & $-0.0012$ \\
세계애니 트렌드 & $-0.0018$ & $+0.0009$ \\
\hline
\end{tabular}
\end{center}

\textbf{노트 498 이 채택한 바로 그 짝이 $1$ 등에서 꼴찌로 간다.} 사전
등록(``애니 트렌드가 두 시기 모두 $1$ 등'')은 \textbf{반증}.
\end{abstract}

\begin{figure}[t]
\centering
\includegraphics[width=\textwidth]{figs/f1.pdf}
\caption{\textbf{가} --- 두 시기. \textbf{나} --- 순위 이동.
\textbf{다} --- 애니 트렌드 세 번 잰 값.}
\end{figure}

\section{채택은 서고, 채택의 \emph{이야기}는 안 선다}

노트 498 은 \texttt{DOMDROP = \{"애니": ("trend\_",)\}} 를 채택했고 그
근거는 두 층이었다.

\begin{enumerate}
\item \textbf{관문} --- 유보에서 자 다섯이 전부 양수(판 $+0.0045$ ·
날짜 통제 $+0.0045$ · 무리 안 $+0.0088$ · KR 만화 $+0.0152$ · 비게임 앱
$+0.0039$).
\item \textbf{내력} --- ``유보를 안 보고 학습만 보고 골랐다.''
\end{enumerate}

\textbf{①은 그대로 선다.} 관문은 \emph{어떻게 고랐는지와 무관하게}
유보에서 잰 값이고, 짝 씨앗 $10$ 에 부호 $10/10$ 이다. \textbf{되돌리지
않는다.}

\textbf{②는 고쳐 적어야 한다.} 옳은 문장은 ``학습만 보고 골랐다''가
아니라 \textbf{``학습 안의 특정 시기 하나가 골랐다''}이다. 시기를 두 해
앞으로 옮겼으면 이 짝은 꼴찌라서 안 뽑혔고, 대신 모바일 media\_push 가
뽑혔을 것이다.

\begin{quote}
\textbf{조항} --- \textbf{안쪽 검증으로 고른 것은 그 시기에 딸린 값이다.}
다른 시기로 한 번 더 돌려 순위가 서는지 보고 나서야 ``학습만 보고
골랐다''고 적는다. 순위가 안 서면 \emph{고르개가 있다}고 말하지 않는다.
\end{quote}

\section{그럼 왜 $2023\!\sim\!2025$ 가 맞았나}

남는 가설은 하나다 --- \textbf{안쪽 검증은 유보에 붙어 있어야 한다.}
$2023\!\sim\!2025$ 는 유보($2025+$)에 인접하고, $2022\!\sim\!2024$ 는
두 해 떨어져 있다. 맞은 쪽이 인접한 쪽이다.

\textbf{$1/1$ 이다.} 규칙이라 부르지 않는다(노트 494 조항 --- 두 사례로도
규칙이라 안 부르는데 하물며 하나다). 다음 사이클에 \textbf{인접성
가설}로 사전 등록해서 시험한다 --- 안쪽 검증을 $2024\!\sim\!2025$ 한 해로
좁히면 애니 트렌드가 더 세게 나와야 한다.

\section{노트 499 조항을 강화한다}

노트 499 는 ``안쪽 고르개는 \textbf{순위}로만 쓰고 개수는 따로 정한다''고
적었다. 이 노트는 그 절반을 지운다.

\begin{quote}
\textbf{순위로도 못 쓴다.} 두 시기의 순위가 $-0.320$ 으로 거꾸로 가는
자에는 순위가 없다.
\end{quote}

남는 것은 \textbf{개수 하나}다 --- 짝을 하나만 뽑으면 이기고
($+0.0046$), 다섯이면 지고($-0.0078$), 아홉이면 더 진다($-0.0085$).
\textbf{어느 하나를 뽑든 하나만 뽑는 것이 규율이었지, 어느 하나인지를
아는 것이 아니었다.}

\section*{남은 것}

\begin{itemize}
\item \textbf{인접성 가설} --- 안쪽 검증 $2024\!\sim\!2025$ 로 좁혀
사전 등록하고 돌린다. 이게 이 갈래의 마지막 시험이다.
\item \textbf{인접성이 틀리면} 안쪽 검증 갈래를 닫는다. 그때는 노트 497
의 $+0.0161$ 을 가져올 다른 길이 필요하다 --- 유보를 안 보는 길이
아니라, \textbf{유보를 나눠 쓰는} 길(예컨대 $2025$ 상반기로 고르고
하반기로 채점).
\item \textbf{노트 498 의 $n_{\text{val}} \geq 150$ 은 이제 이중으로
의심스럽다} --- 유보 실패를 보고 놓았고, 그 아래 순위 자체가 시기에 딸려
있다.
\end{itemize}

\end{document}
